\section{Exact height transport to the quadratic Chabauty models}
\label{sec:height-normalization}

The fixed minimal elliptic curves and rational points are
\[
 E:y^2=x^3-9x-9,\quad P=(-2,1),\qquad
 E':y^2=x^3-189x+999,\quad P'=(6,9).
\]
The monic even-sextic chart of the genus-two curve and its standard
elliptic quotients are
\begin{align}
 C&:W^2=z^6-27z^4+99z^2-9,\\
 F_1&:Y^2=X^3-27X^2+99X-9,\\
 F_2&:Y^2=X^3+99X^2+243X+81.
 \label{ht-standard-models}
\end{align}
The isomorphisms from the minimal models are
\begin{equation}
 i_1(x,y)=(4x+9,8y),\qquad
 i_2(x,y)=(4x-33,8y).
 \label{ht-isomorphisms}
\end{equation}
Thus both have scale \(u=2\), and translation \(r=9\) or \(-33\),
in the notation \(X=u^2x+r,\ Y=u^3y\).
The two displayed short models are globally minimal: their discriminants
are \(2^4 3^6\) and \(2^4 3^{10}\), respectively. At each prime the
valuation is below twelve, whereas replacing an integral model by one
with smaller integral discriminant would decrease that valuation by
a positive multiple of twelve. The curve substitutions in
\eqref{ht-isomorphisms} are exact polynomial identities.

\paragraph{The height convention.}
We use the cyclotomic character
\[
 \chi_5=\log_5,\quad \log_5(5)=0,\qquad
 \chi_\ell(a)=-v_\ell(a)\log_5(\ell)\quad(\ell\ne5),
\]
the unit-root splitting at five, and the diagonal canonical height
\(H_E(Q)\). This is twice Silverman's local-height normalization,
as explicitly fixed in
\cite[\S8.3, printed p.~35]{BalakrishnanDograQCI}.
It differs from the introductory scalar \(h_p\) of
\cite{MazurSteinTateHeight}: that scalar uses character \(\log_p/p\)
and is \(-1/2\) of the corresponding diagonal pairing.
For an infinite-order rational point \(Q\) and a positive integer \(m\),
after this conversion formula~(1.1) of that source reads
\begin{equation}
 H_E(Q)=-\frac{2}{m^2}
       \log_5\!\left(\frac{\sigma_E(mQ)}{d(mQ)}\right),
 \label{ht-sigma-height}
\end{equation}
where \(mQ\) lies in the formal group at five and in the identity
components at the bad primes, and \(d(mQ)\) is the positive square
root of its minimal-model \(x\)-denominator.

\begin{theorem}[Scalar transport and the differential basis]
\label{ht-transport}
For either isomorphism \(i:E_i\longrightarrow F_i\), let
\(\omega=dx/(2y)\), \(\eta=x\omega\), and let \(s\) be the slope
for which \(\eta+s\omega\) spans the unit-root line. The exact
transport laws are
\begin{align}
 H_{F_i}(iQ)&=H_{E_i}(Q),&
 \log_{F_i}(iQ)&=u^{-1}\log_{E_i}(Q),\\
 s_{F_i}&=u^2s_{E_i}-r,&
 \alpha_{F_i}&=u^2\alpha_{E_i},
 \label{ht-transport-laws}
\end{align}
where \(\alpha_E=H_E(P)/\log_E(P)^2\) for a nonzero infinite-order
rational point. In particular, the two height/logarithm-square ratios
in the standard quotient models are four times those on the minimal
models.
\end{theorem}
\begin{proof}
Direct pullback gives
\begin{equation}
 i^*\omega_F=u^{-1}\omega_E,\qquad
 i^*\eta_F=u\eta_E+(r/u)\omega_E.
 \label{ht-differential-pullback}
\end{equation}
Frobenius commutes with the isomorphism, so uniqueness of the ordinary
unit-root line implies
\[
 i^*(\eta_F+s_F\omega_F)=u(\eta_E+s_E\omega_E).
\]
This is exactly the asserted formula for \(s_F\). Integration from
the origin gives the logarithm formula on the formal group;
multiplication by a common integer extends it to all local points.

The scalar height is intrinsic. In this case its invariance can also
be checked directly in the sigma construction.
Theorem~1.3 of \cite{MazurSteinTateHeight} characterizes the normalized
odd integral sigma function by its differential equation. The changes
of formal parameter are integral at five because \(u=2\) is a unit.
Together with \eqref{ht-differential-pullback} and the slope formula,
the defining differential equation and its leading coefficient give
\begin{equation}
 \sigma_F(iQ)=u^{-1}\sigma_E(Q).
 \label{ht-sigma-transport}
\end{equation}
Indeed the parameter change is odd with leading coefficient \(1/u\),
and both sides have coefficient one when expressed in the parameter
on \(F\). Uniqueness proves the equality.
In the four-point sigma quotient defining the local pairing,
\cite[\S2.3]{MazurSteinTateHeight}, there are two factors in each of
the numerator and denominator; all four factors \(1/u\) cancel.
At primes away from five the construction uses the intrinsic minimal
N\'eron model, \cite[\S2.4]{MazurSteinTateHeight}, and is unchanged
under this isomorphism. Bilinearity extends the equality from the
four-point representations to the rational group. Equivalently all
auxiliary local data may be transported through \(i\).
This proves invariance of \(H\); division by the squared logarithm
proves the formula for \(\alpha\).
The rank-one quadraticity argument requires no assertion that the
chosen point generates the full Mordell--Weil group.
\end{proof}

For local integration, the complementary differential scales by \(u\)
and \(\omega\) by \(1/u\), consistently transporting their cup-dual
normalization. Individual tangential-basepoint constants must still be
fixed explicitly. The theorem does not identify an arbitrary local
software output with an intrinsically normalized scalar.

\begin{proposition}[The pinned raw-output transformation]
\label{ht-pari-inverse}
In the official PARI/GP 2.15.4 source, let \([a,b]\) denote the vector
returned on the minimal model and \([A,B]\) the vector returned on
\(F\) for the corresponding point. For the changes above, the source
performs
\begin{equation}
 A=\frac{a+rb}{u},\qquad B=ub.
 \label{ht-raw-transform}
\end{equation}
Its inverse and the recovery of the intrinsically normalized scalar are
\begin{align}
 a&=uA-rB/u,& b&=B/u,\\
 H_E&=uA-\frac{r+s_E}{u}B
     =uA-\frac{(u^2+1)r+s_F}{u^3}B.
 \label{ht-recovered-scalar}
\end{align}
Consequently the uncorrected expression \(A-s_FB\) cannot in general
be used as the canonical scalar on this nonminimal model.
\end{proposition}
\begin{proof}
This first assertion is a version-specific source statement.
In \cite{PARIHeightSource2154},
\texttt{src/\allowbreak basemath/\allowbreak ellpadic.c},
lines 688--689 obtain the global
minimal model and transform the point to it. Lines 741--748 perform
the final transformation \eqref{ht-raw-transform}.
The direction of the change is also explicit in
\texttt{src/\allowbreak basemath/\allowbreak elliptic.c},
lines 4993--5016, and in the official
\texttt{ellchangecurve} documentation \cite{PARIEllipticHeights}:
the coordinates satisfy \(X=u^2x+r\),
\(Y=u^3y+su^2x+t\). Here the latter \(s,t\) are zero and are
unrelated to the unit-root slope.
Inverting the two-by-two transformation is ordinary rational linear
algebra. On the minimal model the canonical scalar is \(a-s_Eb\);
substitution proves the first recovery formula. Substituting
\(s_E=(s_F+r)/u^2\) proves the second.

For the scalar convention, on the minimal short model the source
computes
\[
 a=-2m^{-2}\log_5(\sigma_0(mP)/d),\qquad
 b=-m^{-2}\log_E(mP)^2.
\]
The routine \texttt{ellformallogsigma\_t}, lines 427--440 of
\texttt{ellpadic.c}, implements the zero-slope differential equation.
Replacing \(\sigma_0\) by
\(\sigma_s=\sigma_0\exp(-s_E\log_E^2/2)\) gives \(a-s_Eb\)
and exactly \eqref{ht-sigma-height}.
Thus the raw final vector transformation and the unit-root slope
do not yield the same scalar through the uncorrected combination.
The discrepancy is not a rounding effect. This conclusion is about
these pinned source bytes and inputs; it is not a statement about
every PARI release or every convention for a vector-valued height.
\end{proof}

The simplest implementation computes \([a,b]\) and \(s_E\) on the
minimal model, forms \(H_E=a-s_Eb\) there, and then uses
\eqref{ht-transport-laws} for the standard-model logarithm and ratio.
It need not compute the nonminimal raw vector at all.

\begin{theorem}[An exact leading-digit discrepancy]
\label{ht-mod25}
For each of the actual points \(P,P'\), the difference
\(A-s_FB-H_E(P)\) has five-adic valuation exactly one. This statement
has a finite integer proof independent of PARI's height algorithm.
\end{theorem}
\begin{proof}
Write the exact ninefold point as
\(9P=(A_0/d^2,B_0/d^3)\), with \(d>0\), and similarly for \(P'\).
Integer elliptic addition and reduction modulo 25 give
\[
\begin{array}{c|cc}
 &E&E'\\ \hline
 -A_0/B_0\pmod {25}&14&4\\
 \log_5(-A_0/B_0)\pmod {25}&10&20\\
 H(P)\pmod {25}&5&10
\end{array}
\]
We justify the last row with an error bound for the entire omitted
sigma series. The exact ninefold coordinates are nonsingular at both
bad primes 2 and 3: if a prime divides the denominator the point
reduces to the origin; otherwise at least one of
\(2y\), \(3x^2+a\) is nonzero in the residue field. These alternatives
are checked by integer arithmetic in the certificate below.
The formal parameter \(t\) at five has valuation one, as proved in
the preceding closure calculation. The normalized sigma function is
odd and belongs to \(t\mathbb Z_5[[t]]\), by
\cite[Theorem~1.3]{MazurSteinTateHeight}. Therefore
\[
 \frac{\sigma(t)}t\in1+t^2\mathbb Z_5[[t]],\qquad
 \log_5\frac{\sigma(t)}t\in25\mathbb Z_5.
\]
Since \(t/d=-A_0/B_0\), formula \eqref{ht-sigma-height} gives
\[
 H(P)\equiv-\frac2{81}\log_5(-A_0/B_0)\pmod {25}.
\]
For any unit \(q\in\mathbb Z_5^\times\),
\(\log_5(q)\equiv(q^4-1)/4\pmod {25}\): write \(q^4=1+5k\)
and bound every subsequent logarithm term in \(25\mathbb Z_5\).
This evaluates all entries using modular inverses. The same argument
applies to \(P'\).

Finally \(b=-\log_E(P)^2\) has valuation two, both slopes are
integral, and \(u=2\) is a unit. Equations
\eqref{ht-raw-transform}--\eqref{ht-recovered-scalar} imply
\[
 A-s_FB-H_E(P)\equiv(1/2-1)H_E(P)\pmod {25}.
\]
The residues are 10 and 20, respectively. Each is a nonzero multiple
of five modulo 25, proving the exact valuation. No comparison between
two numerical truncations is used for this conclusion.
\end{proof}

\paragraph{Reproducible evidence and remaining scope.}
The standard-library script
\texttt{replay\_height\_\allowbreak transport.py --check}
first performs an independent binary elliptic addition, the exact
modulo-25 proof, and rational inverse-matrix checks. It then checks
the source-predicted relations with explicit PARI five-adic balls at
precisions 12 and 24. In each run every relation must vanish to the
common declared absolute precision; mere agreement of the two runs
is not the proof of Theorem~\ref{ht-mod25}. The canonical certificate
has SHA256
\texttt{70ec75594f205336\allowbreak 757496605750b130\allowbreak
2e5f3b271afd6acf\allowbreak 5b12275a36c9de05}.
The high-precision sigma and Frobenius routines remain trusted PARI
computations. The exact algebra, scalar transport and leading-digit
error proof are the ordinary mathematical arguments given above.
No new Lean verification of height theory is claimed.

The official source archive used for this analysis has SHA256
\texttt{c3545bfee0c6dfb4\allowbreak 0b77fb4bbabaf999\allowbreak
d82e60069b9f6d28\allowbreak bcb6cf004c8c5c0f}.
Local copies of the relevant C files and function documentation were
inspected during the audit. They can be recovered from the pinned
versioned archive; those copies are not distributed with this publication.
No installed library was patched.
The result supplies the global height coefficients in the prospective
quadratic Chabauty expression. Local-height value sets, compatible
local integrals and chart constants require their own proofs;
certification of every analytic zero and the rational sieve remain
further tasks. Neither a Jacobian rank calculation nor a complete
rational-point classification is asserted here.
